\Article{7}
We begin by evaluating the integral $\int\mathrm{d}s.\varphi(\mu,\mathrm{d}s)$ over all parts of the region $s$, with $\mu$ denoting a given point outside or inside the region $s$. Imagine a spherical surface of radius $=1$ de\SourceLeaf{44}scribed about the center $\mu$ and divided into infinitely small elements. Let $\mathrm{d}\Pi$ be such an element, and let a straight line drawn from $\mu$ toward a point of this element cut the surface of the region $s$ successively at the points $p',p'',p'''$ etc. Their number will be even or odd according as $\mu$ lies outside or inside the region $s$. We shall denote the distances $\mu p',\mu p'',\mu p'''$ etc. by $r',r'',r'''$ etc. Now draw straight lines from $\mu$ toward each point on the boundary of the element $\mathrm{d}\Pi$. They form a pyramidal region and cut out elements of the surface of the region $s$ at the points $p',p'',p'''$ etc., which we shall denote respectively by $\mathrm{d}t',\mathrm{d}t'',\mathrm{d}t'''$ etc. Finally let $q'$ be the angle between the straight line $p'\mu$ and the outward normal to the element $\mathrm{d}t'$. Likewise, let $q'',q'''$ etc. be the inclinations of the corresponding normals at the points $p'',p'''$ etc. to the straight line drawn toward $\mu$. Clearly, then,
\[
\mathrm{d}\Pi=\pm\frac{\mathrm{d}t'.\cos q'}{r'r'}=\mp\frac{\mathrm{d}t''.\cos q''}{r''r''}=\pm\frac{\mathrm{d}t'''.\cos q'''}{r'''r'''}\ \text{etc.}
\]
where the upper or lower signs apply according as $\mu$ is outside or inside the region $s$.

Further, it is clear that the integral $\int\mathrm{d}s.\varphi(\mu,\mathrm{d}s)$ over the parts of the region $s$ contained within that pyramidal region is given by the integral $\mathrm{d}\Pi.\int rr\varphi r.\mathrm{d}r$ taken from $r=r'$ to $r=r''$, then from $r=r'''$ to $r=r''''$ etc., if $\mu$ lies outside the region $s$, or taken from $r=0$ to $r=r'$, then from $r=r''$ to $r=r'''$ etc., if $\mu$ lies inside the region $s$. If, then, we set the indefinite integral
\[
\int rr\varphi r.\mathrm{d}r=-\psi r
\]
with the constant of integration chosen at will, the integral $\int\mathrm{d}s.\varphi(\mu,\mathrm{d}s)$, insofar as it extends over the parts of the region $s$ within that pyramidal region, will be
\[
\begin{aligned}
&=\mathrm{d}\Pi.(\psi r'-\psi r''+\psi r'''-\ \text{etc.})\\
&=\frac{\mathrm{d}t'.\cos q'.\psi r'}{r'r'}+\frac{\mathrm{d}t''.\cos q''.\psi r''}{r''r''}+\frac{\mathrm{d}t'''.\cos q'''.\psi r'''}{r'''r'''}+\ \text{etc.}
\end{aligned}
\]
whenever $\mu$ lies outside the region $s$, and
\[
\begin{aligned}
&=\mathrm{d}\Pi.(\psi0-\psi r'+\psi r''-\psi r'''+\ \text{etc.})\\
&=\mathrm{d}\Pi.\psi0+\frac{\mathrm{d}t'.\cos q'.\psi r'}{r'r'}+\frac{\mathrm{d}t''.\cos q''.\psi r''}{r''r''}+\frac{\mathrm{d}t'''.\cos q'''.\psi r'''}{r'''r'''}+\ \text{etc.}
\end{aligned}
\]
whenever $\mu$ lies inside the region $s$.

Now, on summing over all parts of the spherical surface, the complete integral $\int\mathrm{d}s.\varphi(\mu,\mathrm{d}s)$ becomes
\SourceLeaf{45}
\[
\begin{aligned}
\text{in the first case}\qquad&=\int\frac{\mathrm{d}t.\cos q.\psi r}{rr}\\
\text{in the second case}\qquad&=4\pi\psi0+\int\frac{\mathrm{d}t.\cos q.\psi r}{rr}
\end{aligned}
\]
where $\mathrm{d}t$ denotes any element of the surface of the region $s$, and $q,r$ denote for these elements the same quantities as the primed letters did earlier for particular elements. Finally, $\pi$ denotes half the circumference of a circle of radius $=1$.

It is also readily seen that, if the point $\mu$ lies neither outside the region $s$ nor inside it, but on its surface, the second formula holds with the factor $4\pi$ changed to $2\pi$, provided the surface at $\mu$ has neither a cusp nor an edge. For our purpose there is no need to consider this case.

\Article{8}
By the preceding article, the evaluation of the integral $\int\int\mathrm{d}s.\mathrm{d}S.\varphi(\mathrm{d}s,\mathrm{d}S)$ reduces to
\[
4\pi\sigma\psi0+\int\int\mathrm{d}t.\mathrm{d}S.\frac{\cos q.\psi(\mathrm{d}t,\mathrm{d}S)}{(\mathrm{d}t,\mathrm{d}S)^2}
\]
if we denote by $\sigma$ the volume of the region common to both regions $s,S$. The first part $4\pi\sigma\psi0$ thus drops out if the regions $s,S$ exclude one another. The remaining integral is new and still appears to be double, but it is in fact fivefold. To reduce it to a fourfold integral we shall consider the integral
\[
\int\mathrm{d}S.\frac{\cos q.\psi(\mu,\mathrm{d}S)}{(\mu,\mathrm{d}S)^2}
\]
over all elements of the region $S$. Again, $\mu$ denotes a given point, and $q$ the angle between two straight lines proceeding from it, one toward the element $\mathrm{d}S$ and the other fixed. We shall now show, in two entirely different ways, how to reduce this integral, which appears single but is actually triple, to an integral that is genuinely double.

Let the plane through the point $\mu$ perpendicular to that fixed straight line be denoted by $\Pi$. Imagine the portion of this plane reached by the projection of the region $S$ divided into infinitely small elements $\mathrm{d}\Pi$. Through a point of such an element $\mathrm{d}\Pi$ draw a straight line perpendicular to the plane $\Pi$. Let it cut the surface of the region $S$ successively, that is, as we proceed in a direction parallel to the fixed straight line, at the points $P',P'',P'''$ etc., whose distances from the point $\mu$ are respectively $R',R'',R'''$ etc. Similar straight lines through all points of the boundary of the ele\SourceLeaf{46}ment $\mathrm{d}\Pi$, drawn at right angles to the plane, will form a prismatic region. At the points $P',P'',P'''$ etc. they will cut elements out of the surface of the region $S$, which we denote by $\mathrm{d}T',\mathrm{d}T'',\mathrm{d}T'''$ etc. Finally, let $\chi'$ be the angle between two straight lines proceeding from the point $P'$, one the outward normal to the element $\mathrm{d}T'$, the other parallel to the fixed straight line. Let $P'',P'''$ etc. have the corresponding angles $\chi'',\chi'''$ etc. Clearly, then,
\[
\mathrm{d}\Pi=-\mathrm{d}T'.\cos\chi'=+\mathrm{d}T''.\cos\chi''=-\mathrm{d}T'''.\cos\chi'''\ \text{etc.}
\]
Divide the prismatic region into infinitely small elements $\mathrm{d}\Pi.\mathrm{d}z$, where $z$ denotes the distance of an arbitrary point from the plane $\Pi$, taken as positive on the side toward which the fixed line points. If we denote the distance of this same point from $\mu$ by $r$, then $z=r\cos q$, and since $rr-zz$ is constant, $z\mathrm{d}z=r\mathrm{d}r$, or $\mathrm{d}\Pi.\mathrm{d}z.\cos q=\mathrm{d}\Pi.\mathrm{d}r$. It follows that our integral $\int\mathrm{d}S.\frac{\cos q.\psi(\mu,\mathrm{d}S)}{(\mu,\mathrm{d}S)^2}$ over the parts of the region $S$ contained in that prismatic region is obtained from the integral $\mathrm{d}\Pi.\int\frac{\mathrm{d}r.\psi r}{rr}$ by taking it from $r=R'$ to $r=R''$, then from $r=R'''$ to $r=R''''$ etc. If, then, we set the indefinite integral
\[
\int\frac{\mathrm{d}r.\psi r}{rr}=-\vartheta r
\]
with the constant of integration chosen at will, our integral over the parts of the region $S$ within the prismatic region will be
\[
\begin{aligned}
&=\mathrm{d}\Pi.(\vartheta R'-\vartheta R''+\vartheta R'''-\ \text{etc.})\\
&=-\mathrm{d}T'.\cos\chi'.\vartheta R'-\mathrm{d}T''.\cos\chi''.\vartheta R''-\mathrm{d}T'''.\cos\chi'''.\vartheta R'''-\ \text{etc.}
\end{aligned}
\]
Collecting these sums over the prisms corresponding to all the elements $\mathrm{d}\Pi$ clearly accounts for all elements of the surface of the region $S$, and we obtain the complete integral
\[
\int\mathrm{d}S.\frac{\cos q.\psi(\mu,\mathrm{d}S)}{(\mu,\mathrm{d}S)^2}=-\int\mathrm{d}T.\cos\chi.\vartheta R
\]
where $\mathrm{d}T$ denotes any element of the surface of the region $S$, $R$ its distance from the point $\mu$, and $\chi$ the angle between the outward normal to the element $\mathrm{d}T$ and a straight line parallel to the fixed line.

In this way the integral $\int\int\mathrm{d}s.\mathrm{d}S.\varphi(\mathrm{d}s,\mathrm{d}S)$ has been reduced to the form
\[
4\pi\sigma\psi0-\int\int\mathrm{d}t.\mathrm{d}T.\cos\chi.\vartheta(\mathrm{d}t,\mathrm{d}T)
\]
\SourceLeaf{47}where $\chi$ clearly denotes the inclination of the elements $\mathrm{d}t,\mathrm{d}T$ to one another, measured by the inclination of their normals drawn outward with respect to the regions $s,S$, respectively. The integrations must extend over the entire surfaces of both regions.

